\usepackage[T1]{fontenc}
\usepackage[utf8]{inputenc}
\usepackage{lmodern}
\usepackage[
  paperwidth=7in,
  paperheight=10in,
  inner=0.94in,
  outer=0.84in,
  top=0.84in,
  bottom=0.96in,
  headheight=15pt
]{geometry}
\usepackage{amsmath,amssymb,amsthm,mathtools,bm}
\usepackage{booktabs,array,longtable,multirow}
\usepackage{graphicx,float,caption}
\usepackage{microtype,setspace}
\usepackage{xcolor,enumitem,fancyhdr}
\usepackage{listings}
\usepackage{xurl}
\usepackage{hyperref}

\definecolor{navy}{HTML}{0B1F33}
\definecolor{blue}{HTML}{135E96}
\definecolor{teal}{HTML}{087E8B}
\definecolor{gold}{HTML}{D6A84B}
\definecolor{ink}{HTML}{17202A}
\definecolor{mist}{HTML}{E9F0F5}
\definecolor{softgray}{HTML}{F4F6F7}
\definecolor{codegreen}{HTML}{1B6B45}

\hypersetup{
  colorlinks=true,
  linkcolor=blue,
  urlcolor=teal,
  citecolor=blue,
  pdftitle={Quantitative Finance: From Mathematical Foundations to Empirical Research},
  pdfauthor={Badr Ettousy},
  pdfsubject={Financial engineering, quantitative finance, econometrics, and risk},
  pdfkeywords={quantitative finance, stochastic calculus, derivatives, econometrics, portfolio risk}
}

\setstretch{1.11}
\setlength{\emergencystretch}{1.5em}
\setlength{\parindent}{1.25em}
\setlength{\parskip}{0.15em}
\raggedbottom
\setlist[itemize]{leftmargin=1.7em,itemsep=0.2em,topsep=0.35em}
\setlist[enumerate]{leftmargin=1.9em,itemsep=0.25em,topsep=0.35em}

\pagestyle{fancy}
\fancyhf{}
\fancyhead[L]{\small\sffamily\nouppercase{\leftmark}}
\fancyhead[R]{\small\sffamily Financial Engineering Lab}
\fancyfoot[C]{\sffamily\thepage}
\renewcommand{\headrulewidth}{0.35pt}
\fancypagestyle{plain}{\fancyhf{}\fancyfoot[C]{\sffamily\thepage}\renewcommand{\headrulewidth}{0pt}}

\makeatletter
\renewcommand\chapter{\if@openright\cleardoublepage\else\clearpage\fi
  \thispagestyle{plain}\global\@topnum\z@\@afterindentfalse\secdef\@chapter\@schapter}
\def\@makechapterhead#1{\vspace*{18pt}{\parindent\z@\raggedright\normalfont
  \color{gold}\rule{\textwidth}{1.1pt}\par\vspace{11pt}
  \ifnum\c@secnumdepth>\m@ne {\sffamily\bfseries\Large\color{blue}\@chapapp\space\thechapter\par\nobreak\vskip 8pt}\fi
  {\sffamily\bfseries\huge\color{navy}#1\par\nobreak\vskip 25pt}}}
\def\@makeschapterhead#1{\vspace*{18pt}{\parindent\z@\raggedright\normalfont
  \color{gold}\rule{\textwidth}{1.1pt}\par\vspace{11pt}
  {\sffamily\bfseries\huge\color{navy}#1\par\nobreak\vskip 25pt}}}
\makeatother

\theoremstyle{definition}
\newtheorem{definition}{Definition}[chapter]
\newtheorem{assumption}[definition]{Assumption}
\newtheorem{example}[definition]{Worked example}
\newtheorem{exercise}[definition]{Exercise}
\theoremstyle{plain}
\newtheorem{theorem}[definition]{Theorem}
\newtheorem{proposition}[definition]{Proposition}
\newtheorem{lemma}[definition]{Lemma}
\newtheorem{corollary}[definition]{Corollary}
\theoremstyle{remark}
\newtheorem{remark}[definition]{Remark}
\newtheorem{warning}[definition]{Model-risk warning}

\newenvironment{objectives}
  {\begin{quote}\color{navy}\small\sffamily\textbf{Learning objectives.}\ }
  {\end{quote}}
\newenvironment{researchnote}
  {\begin{quote}\color{teal}\small\sffamily\textbf{Research note.}\ }
  {\end{quote}}
\newenvironment{keypoint}
  {\begin{quote}\color{blue}\small\sffamily\textbf{Key point.}\ }
  {\end{quote}}

\lstset{
  basicstyle=\ttfamily\small,
  keywordstyle=\color{blue}\bfseries,
  commentstyle=\color{codegreen},
  stringstyle=\color{teal},
  backgroundcolor=\color{softgray},
  frame=single,
  rulecolor=\color{mist},
  breaklines=true,
  showstringspaces=false,
  columns=fullflexible,
  keepspaces=true,
  aboveskip=0.8em,
  belowskip=0.8em
}

\DeclareMathOperator{\E}{\mathbb E}
\DeclareMathOperator{\Var}{Var}
\DeclareMathOperator{\Cov}{Cov}
\DeclareMathOperator{\Corr}{Corr}
\newcommand{\Prob}{\mathbb P}
\DeclareMathOperator{\tr}{tr}
\DeclareMathOperator{\rank}{rank}
\DeclareMathOperator{\diag}{diag}
\DeclareMathOperator{\argmin}{arg\,min}
\DeclareMathOperator{\argmax}{arg\,max}
\DeclareMathOperator{\esssup}{ess\,sup}
\DeclareMathOperator{\VaR}{VaR}
\DeclareMathOperator{\ES}{ES}
\DeclareMathOperator{\sgn}{sgn}
\DeclareMathOperator{\logit}{logit}
\DeclareMathOperator{\plim}{plim}
\newcommand{\R}{\mathbb R}
\newcommand{\N}{\mathbb N}
\newcommand{\Q}{\mathbb Q}
\newcommand{\F}{\mathcal F}
\newcommand{\B}{\mathcal B}
\newcommand{\dd}{\,\mathrm d}
\newcommand{\ind}{\mathbf 1}
\newcommand{\ones}{\bm 1}
\newcommand{\transpose}{\mathsf T}
\newcommand{\e}{\mathrm e}
\newcommand{\given}{\,\middle|\,}
\newcommand{\norm}[1]{\left\lVert#1\right\rVert}
\newcommand{\abs}[1]{\left|#1\right|}
\newcommand{\inner}[2]{\left\langle#1,#2\right\rangle}
\newcommand{\Rplus}{\mathbb R_{+}}
\newcommand{\riskfree}{r_f}
